\section{Experimental Setup}

\subsection{Datasets}

Three publicly accessible sources are used, summarised in
Table~\ref{tab:datasets}. All are downloadable without credentialing.

The \textbf{Roboflow DFU} corpus \cite{roboflow_dfu} supplies 9{,}881 labelled images and is
the subject of the derivation study. Content-based grouping reduces it to
3{,}614 source photographs, a property not visible from the file listing
and central to the leakage analysis. It is available at
\url{https://universe.roboflow.com/ssk-r6ppk/diabetic-ulcers}.

The \textbf{Kaggle DFU} dataset \cite{alzubaidi2020} contributes 543
lesion-free skin patches at $224 \times 224$ pixels, used exclusively as
the negative control. No image from this source enters training.

\textbf{DFUC 2024} \cite{dfuc2024} provides the expert-annotated infection task used for
the control experiment, accessed at
\url{https://universe.roboflow.com/wed-axnxx/dfuc-2024}. Its 5{,}372
infection-axis images resolve to 4{,}542 evaluation units from 1{,}417
patients after conflict removal, with 2{,}767 infected and 2{,}605 not
infected, a ratio of $1.06:1$.

Content comparison between the two wound corpora finds 14 shared
photographs, corresponding to 0.2\% of Roboflow and 0.3\% of DFUC, so the
control task is effectively independent of the derivation study.

\begin{table}[!t]
\caption{Dataset summary after content-based deduplication}
\label{tab:datasets}
\centering
\small
\begin{tabular}{@{}llrr@{}}
\toprule
Dataset & Role & Files & Photographs \\
\midrule
Kaggle DFU   & Negative control    & 543      & --      \\
Roboflow DFU & Severity derivation & 9{,}881  & 3{,}614 \\
DFUC 2024    & Control experiment  & 5{,}372  & 4{,}542 \\
\bottomrule
\end{tabular}
\end{table}

Consolidation to one label per photograph yields 30 mild, 1{,}235 moderate
and 2{,}349 severe cases (Fig.~\ref{fig:classes}). The mild class amounts
to approximately six
photographs per fold, which is discussed in Section IX.

\subsection{Implementation and Hardware}

Experiments ran on Google Colab with an NVIDIA Tesla T4 (16\,GB) and, for
non-accelerated stages, on an Apple Silicon workstation using the Metal
Performance Shaders backend. Software: Python 3.10, PyTorch 2.x
\cite{paszke2019}, torchvision, scikit-learn \cite{pedregosa2011},
scikit-image, imbalanced-learn. Hyperparameters are listed in
Table~\ref{tab:hyper}. Mixed precision is enabled on CUDA only.

\begin{table}[!t]
\caption{Training configuration}
\label{tab:hyper}
\centering
\small
\begin{tabular}{@{}ll@{}}
\toprule
Setting & Value \\
\midrule
Backbone            & EfficientNet-B0, ImageNet weights \\
Backbone state      & frozen (unfrozen in Section VII-D) \\
Input size          & $224 \times 224$ RGB \\
Optimiser           & Adam ($10^{-3}$), AdamW ($3\times10^{-5}$) \\
Weight decay        & $10^{-4}$ \\
Batch size          & 256 (probe), 32 (fine-tune) \\
Max epochs          & 60 (probe), 8 (fine-tune) \\
Early stopping      & patience 12 (probe), 3 (fine-tune) \\
Dropout             & 0.3 \\
Tissue encoder      & $3 \to 32 \to 64$, GELU, LayerNorm \\
Gate init $\gamma$  & 0.0 \\
Trainable params    & 86{,}915 ordinal / 86{,}914 binary \\
Cross-validation    & 5-fold, patient-grouped \\
Random seed         & 42 \\
$k$-means           & $k=3$, 4 initialisations, LAB \\
\bottomrule
\end{tabular}
\end{table}

\subsection{Evaluation Metrics}

For the ordinal severity task we report quadratic weighted kappa (QWK),
which penalises distant misgradings more than adjacent ones and is standard
for ordinal clinical scales; macro F1; mean absolute error in grade units;
and a collapsed severe-versus-rest F1, which is the only per-class quantity
the mild sample size supports.

For the binary infection task we report AUROC and AUPRC, which are
threshold-free; Matthews correlation coefficient, which uses all four
confusion-matrix cells and cannot be inflated by the majority class;
sensitivity and specificity separately, because missing an infection and
over-calling one carry different clinical costs; and balanced accuracy.
Every metric is computed on photograph-level scores.

Confidence intervals are percentile bootstrap over folds with 2{,}000
resamples. Paired comparisons across the shared folds use a paired
$t$-test, and the family of tissue-pathway tests is corrected by the
Holm-Bonferroni procedure \cite{holm1979}.

\subsection{Code Availability}

Source code is not publicly available at time of submission. The pipeline
is implemented as eleven sequential notebooks, each asserting its inputs
and writing artefacts consumed by the next, so that every reported value
traces to a specific executed cell.

\section{Results}

\subsection{Label Validation}

Table~\ref{tab:validation} reports the four checks and Fig.~\ref{fig:valid}
summarises three of them. All four fail.

\begin{table}[!t]
\caption{Label validation. Four independent checks, all failing.}
\label{tab:validation}
\centering
\small
\begin{tabular}{@{}llr@{}}
\toprule
Check & Measure & Result \\
\midrule
Expert agreement & derived vs.\ reviewer A & $-0.0254$ \\
Expert agreement & derived vs.\ reviewer B & $+0.0359$ \\
Expert agreement & reviewer A vs.\ B       & $+0.2537$ \\
\midrule
Negative control & lesion-free graded severe & 60.2\% \\
Negative control & max necrosis, intact skin & 0.570 \\
\midrule
Stability & photographs with $>1$ grade & 61.7\% \\
\midrule
Plausibility & mean necrotic fraction     & 0.231 \\
Plausibility & images above severe cut    & 63.1\% \\
\bottomrule
\end{tabular}
\end{table}

\emph{Expert agreement.} On 134 matched images, agreement between the
derived grade and reviewer A was $\kappa = -0.025$, and against reviewer B
$\kappa = +0.036$. Both are indistinguishable from chance. The two
reviewers agreed with each other at $\kappa = 0.254$, which is fair on the
Landis and Koch scale \cite{landis1977}. That value is the relevant
ceiling: severity grading from a photograph is a genuinely difficult task
on which trained clinicians agree only moderately, and the derived labels
do not approach even that level.

\emph{Negative control.} Applied to 543 lesion-free skin images, the
derivation returned a severe grade for 60.2\% of them, with a mean necrotic
fraction of 0.222 and a maximum of 0.570 on intact tissue. The correct
value in every case is zero. These images are out of domain by design; the
purpose of a negative control is not to measure performance but to observe
behaviour where ground truth is known. The same contamination is present in
wound images, where it cannot be detected because no reference exists.

\emph{Stability.} Of 2{,}818 photographs with more than one augmented copy,
1{,}739 (61.7\%) received more than one severity grade across copies of
themselves. Augmentation does not change the wound, so a label that changes
under augmentation is not measuring the wound.

\emph{Plausibility.} The corpus mean necrotic fraction is 0.231 and 63.1\%
of images exceed the severe threshold. No clinical DFU cohort is
approximately two-thirds severe.

\subsection{Mechanism}

Three mechanisms explain these results, each sufficient alone.

First, $k$-means with $k = 3$ partitions every image into three non-empty
clusters whose proportions sum to one. There is no output corresponding to
``no necrosis present''. The darkest third of the pixel distribution is
labelled necrotic by construction, whatever the image contains. This is why
lesion-free skin receives a severe grade.

Second, the assumed luminance ordering is inverted. Cross-referencing
against tissue references places granulation near $L^{*} = 118$ and
periwound callus near $L^{*} = 168$: granulation is darker than callus,
the opposite of the ordering the derivation assumes when it assigns the
highest-luminance cluster to granulation.

Third, all dark regions fall into the same cluster: shadow, wound depth,
hair, image background and darker skin. Fitzpatrick phototypes V and VI
\cite{fitzpatrick1988} overlap the necrosis reference range in LAB space,
so the method over-grades darker skin systematically. This is a fairness
concern as well as a validity one, and it would not be visible in aggregate
accuracy on a predominantly light-skinned cohort.

\subsection{Control Experiment}

The results above show that the derived labels do not recover severity.
They do not by themselves show that the labels are the only problem. To
separate label quality from pipeline quality, the identical architecture,
preprocessing and evaluation protocol were applied to the DFUC 2024
infection task, which carries clinician annotation.

Infection and severity are different constructs and this paper does not
treat them as interchangeable. Infection is one of several criteria
contributing to a severity grade, alongside depth, ischaemia and probing to
bone. The infection task is used here solely as a control condition: it
supplies expert-assigned labels on wound photographs of the same type and
resolution, which is the property required to test whether the pipeline
functions when labels are valid. A model that detects infection does not
grade severity, and no result in this section should be read as a severity
grading result.
Table~\ref{tab:control} reports the comparison and Fig.~\ref{fig:control}
presents it graphically.

\begin{table}[!t]
\caption{Control comparison. Identical pipeline, identical folds, identical
code. Only the label source differs.}
\label{tab:control}
\centering
\small
\begin{tabular}{@{}llr@{}}
\toprule
Task & Metric & Value \\
\midrule
Severity (derived colour labels)  & QWK                 & 0.2492 \\
Severity (derived colour labels)  & severe-vs-rest F1   & 0.6986 \\
\midrule
Infection (expert labels)         & AUROC               & 0.8101 \\
Infection (expert labels)         & AUPRC               & 0.8685 \\
Infection (expert labels)         & MCC                 & 0.4637 \\
% FLAG: sensitivity/specificity below are from the SUPERSEDED run.
% Rerun In\cite{anilkumar2026} / notebook 09 and replace before submission.
Infection (expert labels)         & sensitivity         & 0.7083\% $^{\dagger}$ \\
Infection (expert labels)         & specificity         & 0.7736\% $^{\dagger}$ \\
\bottomrule
\end{tabular}
\end{table}

From pixels alone, the same frozen representation and the same linear head
reach QWK $0.249$ on derived severity labels and AUROC $0.810$
(95\% CI $0.797$--$0.829$, bootstrap over folds) on expert infection
labels.

Because the architectures, preprocessing, grouping and code are unchanged
between these two rows, the difference isolates label provenance. Model
capacity, input resolution, preprocessing and training configuration are
ruled out as explanations.

\subsection{Effect of Unfreezing the Backbone}

All results above use a frozen backbone, making them linear probes on
ImageNet features and therefore a floor rather than a tuned ceiling. To
test whether that floor is the binding constraint, the full backbone was
fine-tuned end to end on the same folds with the same per-photograph
scoring.

Fine-tuning reached AUROC $0.8180$, AUPRC $0.8638$, MCC $0.4869$,
sensitivity $0.7414$ and specificity $0.7546$,
% FLAG: these fine-tuning figures are from the SUPERSEDED frozen baseline
% run. The fine-tuned model itself may not need rerunning, but confirm
% 0.8180 etc. are still current before relying on the delta below.
an improvement of $+0.0079$
AUROC over the frozen probe (recomputed against the current baseline of
$0.8101$). Training 5.29 million parameters rather than approximately 87
thousand yields an improvement smaller than the per-fold standard
deviation reported for the frozen model.
% FLAG: "smaller than the per-fold standard deviation" needs the current
% per-fold SD (see flag above) to confirm this claim still holds.
The representation is not the
bottleneck.

\subsection{Ablation}

Table~\ref{tab:ablation} reports all nine configurations, six on the
severity task and three on the control task. Every row is an executed run.

\begin{table}[!t]
\caption{Ablation. Every row is an executed run. The flagged row is
circular and is discussed in the text.}
\label{tab:ablation}
\centering
\small
\begin{tabular}{@{}llr@{}}
\toprule
Task & Configuration & Metric \\
\midrule
\multicolumn{3}{@{}l}{\emph{Severity (derived labels), QWK}} \\
& CNN only (softmax)          & 0.2492 \\
& CNN only (CORAL)            & 0.0000 \\
& + tissue, concat            & 0.9458$^{\dagger}$ \\
& + tissue, gated (TGO-Net)   & 0.8955$^{\dagger}$ \\
& + tissue, gated + SMOTE     & 0.9075$^{\dagger}$ \\
& CNN only + SMOTE            & 0.1290 \\
\midrule
\multicolumn{3}{@{}l}{\emph{Infection (expert labels), AUROC}} \\
& CNN only                    & \textbf{0.8101} \\
& + tissue, concat            & 0.8101 \\
& + tissue, gated (TGO-Net)   & 0.7945 \\
\bottomrule
\multicolumn{3}{@{}l}{\footnotesize $^{\dagger}$ circular: see
Section VII-E.}
\end{tabular}
\end{table}

On the severity task, adding the tissue proportions raises QWK from
$0.0000$ to $0.9458$. This is not evidence that tissue features improve
severity classification. The tissue proportions are the label-generating
variable: (\ref{eq:rule}) defines severity as a threshold on
$\rho_{\text{nec}}$, and mutual information between the necrotic fraction
and the assigned label is $0.565$ nats against approximately $0.10$ nats
for the other two channels. Supplying the model with the quantity that
produced the label yields near-perfect fit by construction. The correct
reading is that this result confirms the labels reduce to a deterministic
threshold on one continuous variable, which is itself part of the negative
finding.

The control task carries no such circularity, because the labels were
assigned by clinicians and not derived from the tissue proportions. There,
the ordering reverses. Adding tissue by concatenation changes AUROC by
$0.0000$, no detectable effect.
% FLAG: the concat p-value is not printed by the statistics cell and is
% still unconfirmed. Obtain it before submission.
Adding tissue through the gate changes AUROC by $-0.0142$ ($p = 0.016$,
paired across the five shared folds), a small but statistically
detectable degradation. The tissue pathway does not help on expert
labels, and on this measure modestly harms performance.

The gate scalar agrees independently. Across folds $\gamma$ has mean
$-0.060$ against a standard deviation of $0.928$, and is not
distinguishable from zero ($p = 0.90$). Fig.~\ref{fig:gate} shows the per-fold values. Given a free
parameter initialised at zero and the opportunity to use the tissue
proportions, the network did not converge on how much to rely on them, and
changed sign on three of four fold transitions. Reporting the mean alone
would misrepresent a parameter that never settled.

\subsection{Explainability}

Grad-CAM was computed on a held-out fold of the infection task using a
model fine-tuned end to end, since a frozen backbone provides no gradient
path to the convolutional layers. Fig.~\ref{fig:gradcam} shows image,
heatmap and cluster map for representative cases.

Area-normalised activation lift was $1.283$ for the necrotic cluster,
$1.110$ for slough and $0.767$ for granulation. Attention is therefore
mildly concentrated on darker regions relative to their area.

This quantity must be interpreted narrowly. The cluster masks are produced
by the same LAB $k$-means shown in Section VII-A to grade 60.2\% of
lesion-free skin as severe and to invert the luminance ordering of tissue
types. A high alignment score therefore indicates that the model attends to
the same dark regions the clustering labels necrotic, which include shadow,
wound depth and darker skin, and not that it attends to necrotic tissue.
The measurement compares two procedures that share a failure mode. It is
reported as a diagnostic of internal consistency, not as evidence of
clinical correctness. Establishing the latter requires masks from a
segmenter trained on expert tissue annotations.

\subsection{Comparison with Existing Studies}

Table~\ref{tab:related} places this work against recent severity grading
studies. All prior work relies on private hospital datasets with
clinician-assigned grades, which is precisely the resource this study
attempted to substitute for and could not.

\begin{table}[!t]
\caption{Comparison with recent severity grading studies}
\label{tab:related}
\centering
\small
\begin{tabular}{@{}llccc@{}}
\toprule
Study & Data & Classes & Reported & Public \\
\midrule
Zhou et al.\ \cite{zhou2025}       & Private (CN) & 5 & 0.94 AUC & No \\
Abebe et al.\ \cite{abebe2025}     & Private (ET) & 6 & 82\%     & No \\
Anilkumar et al.\ \cite{anilkumar2026} & Private (IN) & 3 & 82\% & No \\
Wang et al.\ \cite{wang2024score}  & Mixed        & 4 & 95.3\%   & Partial \\
This work                          & Public       & 3 & n/a$^{\ddagger}$ & Yes \\
\bottomrule
\multicolumn{5}{@{}l}{\footnotesize $^{\ddagger}$ no severity classifier is
reported; see text.}
\end{tabular}
\end{table}

We deliberately report no severity performance figure for comparison. The
labels required to train such a classifier were not recoverable, and
quoting a number obtained against invalid labels would be misleading.
The figures in Table~\ref{tab:related} are also not directly
comparable to results obtained under the protocol of Section V-D: none of
these studies documents whether augmented copies of the same photograph, or
multiple wound sites from the same patient, were separated across splits.
Given the leakage this study found in a public corpus, this is raised as a
general methodological observation rather than a claim about any specific
study.

\section{Discussion}

The central result is that a widely plausible labelling strategy fails, and
that the failure is detectable cheaply. Each of the four checks in
Section VII-A costs minutes of computation, and any one of them would have
prevented training on invalid labels. The negative control in particular
requires only a set of images known to contain no lesion, and it converts
an unverifiable claim about wound images into a verifiable one.

The mechanism matters more than the score, and Fig.~\ref{fig:valid}
summarises the evidence. Severity in the Wagner and
University of Texas systems depends on wound depth, probing to bone,
infection and ischaemia. A photograph does not determine depth. The failure
is therefore one of identifiability rather than tuning: no clustering
scheme over colour recovers a quantity the image does not encode. This
explains why successive recalibration of the thresholds produced small and
saturating gains rather than convergence toward acceptable agreement.

The control experiment is what converts this from a failed attempt into an
interpretable result, and Fig.~\ref{fig:control} states it in one image. Without it, a reader could reasonably conclude that
the pipeline was inadequate. With it, the conclusion is specific: the same
code, the same architecture and the same class of images produce a usable
classifier when the labels carry clinical meaning, and fail when the labels
are derived from colour.

The TGO-Net results deserve care, and Fig.~\ref{fig:gate} is the reason. On the derived labels the gated model
appears excellent, and a paper reporting only that number would claim a
successful tissue-guided architecture. The control task shows the opposite:
the gate modestly degrades performance, and the gate scalar never
settles. Both observations point the same way, and they are only separable
because the architecture was designed with a readable diagnostic. We
therefore present TGO-Net as an instrument rather than a performance
contribution. A zero-initialised gate is a cheap addition to any
multi-modal design and makes modality reliance directly observable.

Finally, the fine-tuning result bounds the scope of the problem. Training
sixty times more parameters yielded an improvement smaller than the
per-fold standard deviation. Whatever limits performance on this task, it
is not the representation.

\section{Limitations}

The central limitation is the one this paper documents. Colour-based tissue
proportion analysis does not recover severity, and no severity classifier is
reported here because no valid severity labels were obtained. The original
objective of a reproducible three-class severity classifier on public data
remains unmet.

The derivation is assessed as implemented, not exhaustively. A different
clustering scheme, a segmenter trained on annotated tissue masks, or
additional sensing modalities might recover tissue composition where this
method does not. The supportable claim is that this widely used approach
fails these checks, and that the checks are cheap enough to apply before
building on any derived label.

The control task is infection rather than severity. Infection is one
component of the grading criteria and was selected because it carries
expert annotation, so it establishes that the pipeline functions with valid
labels but does not itself constitute severity grading.

The mild class was never measurable. Across 3{,}614 photographs it holds
30, approximately six per fold. Feature-space resampling expands the
training set but cannot create independent evaluation cases, and no
architectural change addresses this. Per-class figures for mild are
reported with Wilson intervals wide enough to preclude any claim.

The expert review comprises 134 matched images assessed by two reviewers.
A larger panel would tighten the inter-rater estimate, which is the
quantity that bounds what any automated method could achieve.

Grad-CAM analysis used 16 held-out images, sufficient to characterise the
area-normalised lift but not to support finer claims about attention
patterns.

\section{Future Work}

The most direct route to the original objective is clinical annotation. A
few hundred images graded by two or more clinicians, with inter-rater
agreement reported, would supply ground truth for supervised severity
grading and simultaneously provide a reference against which any derived
label could be assessed. The protocol in Section V-C indicates the sample
size each check requires.

Replacing colour clustering with a segmenter trained on expert tissue masks
would test whether tissue composition is informative when measured rather
than assumed. The negative result here concerns derived proportions, not
the underlying clinical concept, and public tissue segmentation datasets
now exist for this purpose \cite{dhar2024}.

Additional sensing modalities address the identifiability problem directly.
Wound depth is central to grading and absent from a colour photograph;
depth or thermal imaging would supply it. This is a change of input rather
than of model, which is what the present results indicate is required.

Finally, the validation protocol generalises beyond this application.
Expert agreement, a negative control on lesion-free images, and a stability
test across augmented copies are inexpensive and applicable to any
derived-label pipeline in medical imaging.

\section{Conclusion}

This paper set out to grade diabetic foot ulcer severity from colour wound
images by deriving labels through $k$-means clustering in CIE LAB colour
space, since no public dataset provides severity annotation. The derivation
was implemented, applied to 9{,}881 images, and then validated rather than
assumed correct.

It does not recover severity. Agreement with two clinicians was
$\kappa = -0.025$ and $\kappa = +0.036$ against an inter-rater agreement of
$\kappa = 0.254$. A negative control returned a severe grade for 60.2\% of
lesion-free skin images, assigning necrotic fractions up to 0.570 to intact
tissue. Nearly two photographs in three received more than one grade across
their own augmented copies. The mechanism is identifiable rather than
incidental: $k$-means with $k=3$ cannot report the absence of a tissue
type, the assumed luminance ordering is inverted, and every dark region
including darker skin falls into the cluster labelled necrotic.

Applying the identical pipeline to an expert-annotated infection task gave
AUROC $0.810$ under a verified leakage-free protocol, and unfreezing the
backbone changed this by $+0.008$. TGO-Net, a tissue-guided architecture
with a zero-initialised gate, confirmed the same conclusion from a
different direction: on derived labels it appears to excel because it is
supplied with the label-generating variable, while on expert labels it is
slightly worse than a plain CNN and its gate never settles on a value.

Severity grading from colour wound images is not blocked by architecture,
image resolution or training method. It is blocked by the absence of
clinically annotated ground truth, and label derivation from image content
does not substitute for it. We release the validation protocol that
establishes this, so that the same checks can be applied before other work
builds on a derived label.
